\section{Conclusion}
\label{sec:conclusion}

We investigated whether \aistyle is linearly represented in the residual stream of large language models.
Using contrastive activation analysis on paired human and ChatGPT text, we found that linear probes achieve 93--96\% accuracy at distinguishing the two classes across two architecturally different models, while the one-dimensional difference-in-means direction achieves 69--81\%.
The gap between these numbers suggests that \aistyle occupies a multi-dimensional subspace rather than a single direction, distinguishing it from the cleaner single-direction geometry observed for refusal and truth.

Our results connect mechanistic interpretability to AI-generated text detection.
The strong signal in a base model (\pythia) shows that stylistic representations emerge during pre-training, before alignment, and the dramatically earlier peak in an instruction-tuned model (\qwen) suggests that alignment training reorganizes these representations.

Several directions remain open.
Testing with text from multiple AI generators would clarify whether the direction is model-general or ChatGPT-specific.
Sparse autoencoder decomposition could identify the individual stylistic features (formality, hedging, verbosity) that compose the subspace.
Finally, causal steering on instruction-tuned models would provide a stronger test of whether the direction controls \aistyle in generated text, with potential applications to controllable generation and more robust AI text detection.
